% In this file you should put macros to be used only by
% the web version. Of course they should have a corresponding
% version in macros/print.tex.
% Typically the printed version could have more fancy decorations.
% This will probably be a very short file.

% Inline doc-gen4 cross-reference for the web build. Renders
% \leandocref{Proofs.pdiv\_comp} as a clickable typewriter link to the
% corresponding declaration's API doc page.
\newcommand{\leandocref}[1]{\href{\docbase/find/\#doc/#1}{\texttt{#1}}}

% Cross-reference into the Bestiary chapter. plasTeX renders
% \hyperref[label]{text} as a dead href="#." link, so the web version
% links straight to the anchor on the bestiary chapter page (plasTeX
% names it after the chap:bestiary label; ids keep their colons).
\newcommand{\bestiaryref}[2]{\href{chap-bestiary.html\##1}{#2}}

% \prosesection: chapter-body prose sub-heading. Rendered as a styled bold
% heading that is NOT a sectioning command, so it never produces a web
% sidebar entry -- the ch1-9 web index then shows only numbered sections,
% matching the PDF TOC. The printed version (macros/print.tex) keeps it as
% a normal \section*.
\newcommand{\prosesection}[1]{\par\bigskip\noindent{\large\bfseries #1}\par\nopagebreak\smallskip}

% \proseparagraph: the run-in bold heading of a demo's sub-steps ("Scored per patient.").
% Print makes it a \paragraph; here it is bold text and not a sectioning command, so the
% web index lists chapters, sections and the bestiary's per-network subsections and no
% fifth level -- plasTeX keeps its own \paragraph, a \renewcommand of it changes nothing.
\newcommand{\proseparagraph}[1]{\par\medskip\noindent\textbf{#1}\ }

% \cellleft: see macros/print.tex. HTML table cells are already ragged right.
\newcommand{\cellleft}{}

% \setbookmarksdepth: see macros/print.tex. HTML has no PDF bookmarks.
\newcommand{\setbookmarksdepth}[1]{}
